\chapter*{不等式}
\subsection*{不等式的性质}
\begin{enumerate}
    \item (2024$\cdot$杨浦区二模$\cdot$14)已知实数 $a$ 、 $b$ 、 $c$ 、 $d$ 满足 $a > b > 0 > c > d$ ，则下列不等式一定正确的是\dotfill{(\kh{C})}
    \xx{$a + d > b + c$}{${ad} > {bc}$}{$a + c > b + d$}{${ac} > {bd}$}

    \item (2024$\cdot$崇明区二模$\cdot$13)若 $a > b$ ， $c < 0$ ，则下列不等式成立的是 \dotfill{(\kh{D})}
    \xx{$a > b - c$}{$\dfrac{a}{c} > \dfrac{b}{c}$}{$a + c < b + c$}{$a{c}^{2} > b{c}^{2}$}

    \item (2025$\cdot$嘉定区二模$\cdot$13)已知实数 $a$ 、 $b$ 满足 $a > b$ ，则下列不等式中，不恒成立的是 \dotfill{(\kh{B})}
    \xx{${a}^{3} > {b}^{3}$}{$\dfrac{a}{b} > 1$}{${a}^{2} + {b}^{2} > {2ab}$}{${2}^{a} > {2}^{b}$}

    \item (2024$\cdot$浦东新区一模$\cdot$13)如果 $a > 0 > b$ ，则下列不等式中一定成立的是 \dotfill{(\kh{D})}
    \xx{$\sqrt{a} > \sqrt{-b}$}{${a}^{2} > {b}^{2}$}{${a}^{2} < {ab}$}{${a}^{3} > {b}^{3}$}
\end{enumerate}

\subsection*{解分式不等式}
\begin{enumerate}
    \item 不等式 $\dfrac{5 - x}{{x}^{2} - {2x} - 3} \leq   - 1$ 的解集为\blkda[4]{$\left( {-1,1\rbrack \cup \lbrack 2,3}\right)$}.
    \begin{jieda}
        $\dfrac{5 - x}{{x}^{2} - {2x} - 3} \leq   - 1 \Longleftrightarrow \dfrac{{x}^{2} - {3x} + 2}{{x}^{2} - {2x} - 3} \leq  0 \Longleftrightarrow\left\{  \begin{array}{l} \left( {x + 1}\right) \left( {x - 1}\right) \left( {x - 2}\right) \left( {x - 3}\right)  \leq  0 \\  x \neq   - 1\text{ 且 }x \neq  3 \end{array}\right.$

        穿针引线可得解集为$\left( {-1,1\rbrack \cup \lbrack 2,3}\right)$ 
    \end{jieda} 

    \item 关于 $x$ 的不等式 $\dfrac{\left( {x + 3}\right) {\left( x - 2\right) }^{999}{\left( x + 1\right) }^{2}}{{\left( x - 1\right) }^{2015}{\left( x - 4\right) }^{2016}} \leq  0$ 的解集为\blkda[4]{$(- \infty , - 3 \rbrack  \cup  \{  - 1\}  \cup  (1,2\rbrack$}
    \begin{jieda}
        原不等式等价于: $\left\{  \begin{array}{l} \left( {x + 3}\right) {\left( x - 2\right) }^{999}{\left( x + 1\right) }^{2}{\left( x - 1\right) }^{2015}{\left( x - 4\right) }^{2016} \leq  0, \\  x - 1 \neq  0, \\  x - 4 \neq  0, \end{array}\right. $
        
        结合高次不等式奇穿偶不穿的性质求解不等式组可得原不等式的解集为: $(- \infty , - 3 \rbrack  \cup  \{  - 1\}  \cup  (1,2\rbrack$
    \end{jieda} 
    \item 解关于$x$的不等式：$\dfrac{a}{x-2} > 1-a$
    \begin{jieda}
        原不等式化为$\dfrac{(a-1)x+(2-a)}{x-2} > 0$，进一步化为$[(a-1)x+(2-a)](x-2) > 0$ 
        \begin{dingautolist}{172}
            \item $a > 1$时，不等式等价于$\left[x-\dfrac{a-2}{a-1}\right](x-2) > 0$，又因为$\dfrac{a-2}{a-1} = 1 - \dfrac{1}{a-1} < 1$，所以原不等式的解集为$\left(-\infty, \dfrac{a-2}{a-1}\right) \cup (2, +\infty)$
            \item $a = 1$时，不等式等价于$x-2 > 0$，解集为$(2, +\infty)$
            \item $a < 1$时，不等式等价于$\left[x-\dfrac{a-2}{a-1}\right](x-2) < 0$ \\
            此时$\dfrac{a-2}{a-1} < 2 \Longleftrightarrow a < 0$，因此：\\
            $a < 0$时原不等式解集为$\left(\dfrac{a-2}{a-1}, 2\right)$ \\
            $a = 0$时原不等式解集为$\varnothing$ \\
            $a \in (0, 1)$时原不等式解集为$\left(2, \dfrac{a-2}{a-1}\right)$
        \end{dingautolist}
        \ps 参数会影响二次方程的两个根的大小关系，而最高次项含参则会涉及到退化回一次函数，以及二次函数开口方向.
    \end{jieda}

    \item 已知集合 $A = \left\{  {x\left| {\dfrac{{ax} - 2}{x - a} \leq  0}\right. }\right\}$ ,若 $2 \notin  A$ ,则实数 $a$ 的取值范围是\blkda[4]{$(1,2\rbrack$}.
    \begin{jieda}
        集合 $A = \left\{  {x\left| {\dfrac{{ax} - 2}{x - a} \leq  0}\right. }\right\}$ ,若 $2 \notin  A$ ,故 $\dfrac{{2a} - 2}{2 - a} > 0$ 或 $2 - a = 0$ ,解得 $1 < a \leq  2$ .
    \end{jieda}
\end{enumerate}

\subsection*{已知不等式解集求参数}
\begin{enumerate}
    \item 不等式 $a{x}^{2} + {bx} + 2 > 0$ 的解集是 $\left\{  {x\left| { - \dfrac{1}{2} < x < \dfrac{1}{3}}\right. }\right\}$ ,则 $a + b =$ \blkda[4]{$-14$}.
    \begin{jieda}
        不等式 $a{x}^{2} + {bx} + 2 > 0$ 的解集是 $\left\{  {x\left| { - \dfrac{1}{2} < x < \dfrac{1}{3}}\right. }\right\} \Longleftrightarrow \left\{  \begin{array}{l} a < 0 \\   - \dfrac{b}{a} =  - \dfrac{1}{6}, \\  \dfrac{2}{a} =  - \dfrac{1}{6} \end{array}\right.$ 
        
        解得$\left\{  {\begin{array}{l} a =  - {12} \\  b =  - 2 \end{array}}\right.$，故$a + b =  - {14}$
    \end{jieda}
    \item 若关于 $x$ 的不等式 $a{x}^{2} + {bx} + c < 0$ 的解集是 $\left( {-\infty , - 1}\right)  \cup  \left( {3, + \infty }\right)$ ,则关于 $x$ 的不等式$a{x}^{2} + {cx} + {2b} \geq  0$ 的解集是\blkda[4]{$\left\lbrack  {-1,4}\right\rbrack$}.
    \begin{jieda}
        不等式 $a{x}^{2} + {bx} + c < 0$ 的解集是 $\left( {-\infty , - 1}\right)  \cup  \left( {3, + \infty }\right) \Longleftrightarrow \left\{  \begin{array}{l}  
        a < 0 \\
        a-b+c = 0 \\  
        9+3b+c =  0 
        \end{array}\right.$
        
        解得 $\left\{\begin{array}{l}
             a < 0 \\
             b =  - {2a}  \\
             c =  - {3a} 
        \end{array}\right.$，故$a{x}^{2} + {cx} + {2b} \geq  0 \Longleftrightarrow a{x}^{2} - {3ax} - {4a} \geq  0$
        
        因 $a < 0$ ,故得 ${x}^{2} - {3x} - 4 \leq  0$ ,解得 $x \in \left\lbrack  {-1,4}\right\rbrack$
    \end{jieda}
    \item ${ax} + b < 0$ 的解集为 $\left( {-\infty , - 1}\right)$ ,则 $\left( {{ax} - b}\right) \left( {x + 2}\right)  < 0$ 的解集为\blkda[4]{$(-2,1)$}.
    \begin{jieda}
         ${ax} + b < 0$ 的解集为 $\left( {-\infty , - 1}\right)$，即$\left\{\begin{array}{l}
              b = a  \\
              a > 0
         \end{array}\right.$

         故$\left( {{ax} - a}\right) \left( {x + 2}\right)  < 0 \Longrightarrow  \left( {x - 1}\right) \left( {x + 2}\right)  < 0$，解得$- 2 < x < 1$ 
    \end{jieda}
\end{enumerate}

\subsection*{整根问题}
\begin{enumerate}
    \item 关于 $x$ 的不等式 ${x}^{2} - \left( {1 + {2a}}\right) x + {2a} < 0$ 的解集中恰有 2 个整数,则实数 $a$ 的取值范围是\blkda[4]{$\left\lbrack  {-1, - \dfrac{1}{2}}\right)  \cup  \left( {\dfrac{3}{2},2}\right\rbrack$}.
    \begin{jieda}
        ${x}^{2} - \left( {1 + {2a}}\right) x + {2a} < 0 \Longleftrightarrow \left( {x - 1}\right) \left( {x - {2a}}\right)  < 0$ ,

        \begin{dingautolist}{172}
            \item 当 ${2a} = 1$ 时,不等式化为 ${\left( x - 1\right) }^{2} < 0$ ,此时无解
            \item 当 ${2a} > 1$ 时,解得 $1 < x < {2a}$ ,要使解集中恰有两个整数,则 $3 < {2a} \leq  4$ ,得 $\dfrac{3}{2} < a \leq  2$
            \item 当 ${2a} < 1$ 时,解得 ${2a} < x < 1$ ,要使解集中恰有两个整数,则 $- 2 \leq  {2a} <  - 1$ ,得 $- 1 \leq  a <  - \dfrac{1}{2}$
        \end{dingautolist}
        综上，$a \in \left\lbrack  {-1, - \dfrac{1}{2}}\right)  \cup  \left( {\dfrac{3}{2},2}\right\rbrack$ .
    \end{jieda}
    \item 已知关于 $x$ 的不等式组 $\left\{  \begin{array}{l} {x}^{2} - {2x} - 8 > 0, \\  2{x}^{2} + \left( {{2k} + 7}\right) x + {7k} < 0 \end{array}\right.$ 仅有一个整数解,则实数 $k$ 的取值范围是\blkda[4]{$ [-5, 3) \cup (4,5]$}.
    \begin{jieda}
        不等式 ${x}^{2} - {2x} - 8 > 0$ 的解集为 $\{ x \mid  x <  - 2$ 或 $x > 4\}$ .
        
        不等式 $2{x}^{2} + \left( {{2k} + 7}\right) x + {7k} < 0$ 化为 $\left( {x + k}\right) \left( {x + \dfrac{7}{2}}\right)  < 0$ .

        \begin{dingautolist}{172}
            \item 当 $k = \dfrac{7}{2}$ 时,不等式组的解集为 $\varnothing$ ,不符合题意.
            
            \item 当 $k > \dfrac{7}{2}$ 时,不等式 $2{x}^{2} + \left( {{2k} + 7}\right) x + {7k} < 0$ 的解集为 $\left\{x\left |  { - k < x <  - \dfrac{7}{2}}\right. \right\}$，由解集中的整数为$-4$,得 $- 5 \leq   - k <  - 4$ ,解得 $4 < k \leq  5$ .

            \item 当 $k < \dfrac{7}{2}$ 时,不等式 $2{x}^{2} + \left( {{2k} + 7}\right) x + {7k} < 0$ 的解集为 $\left\{  {x\left| { - \dfrac{7}{2} < x <  - k}\right. }\right\}$，由解集中的整数为$-3$,得 $- 3 <  - k \leq  5$ ,解得 $- 5 \leq  k < 3$ .
        \end{dingautolist}
        综上，$k \in [-5, 3) \cup (4,5]$.
    \end{jieda}

    \item 已知关于$x$的不等式组$\left \{ \begin{array}{l}
         \dfrac{x-a}{x-1+a} < 0\\
         x+3a>1 
    \end{array}\right.$的整数解恰有两个，则实数$a$的取值范围是\blkda[4]{$(1, 2]$}
    \begin{jieda}
        首先可以确定第一个不等式的解集为$(a, 1-a)$或$(1-a, a)$，因为$a$和$1-a$关于$\dfrac{1}{2}$对称，因此$\dfrac{1}{2}$必然在该不等式的解集中。而第二个不等式的解集为$(1-3a, +\infty)$
        \begin{dingautolist}{172}
            \item 若$a \in [0, 1]$，则$(a, 1-a)$或$(1-a, a)$没有整数解，故整个不等式组的解集中也没有整数解.
            \item 若$a < 0$，则$a < 0 < 1 < 1-a < 1-3a$，则整个不等式组解集为空集.
            \item 若$a > 1$，则$1-3a < 1-a < 0 < 1 < a$整个不等式组的解集为$(1-a, a)$，此时解集中必有两个整数解$0, 1$，故当$a \in (1, 2]$时，满足要求.
        \end{dingautolist}
        综上，$a \in (1, 2]$ \\
        \ps 本题画数轴可以更方便分析，本题判断出$a$和$1-a$关于$\dfrac{1}{2}$对称后可以猜测到两个整根是0和1.
    \end{jieda}
\end{enumerate}

\subsection*{解含绝对值不等式}
\begin{enumerate}
    \item 不等式 $\left| {1 - {4x}}\right|  > {2x}$ 的解集为\blkda[4]{$\left(-\infty, \dfrac{1}{6}\right) \cup \left(\dfrac{1}{2}, +\infty\right)$} 
    \begin{jieda}
        $\left| {1 - {4x}}\right|  > {2x} \Longleftrightarrow 1 - {4x} > {2x}$ 或 $1 - {4x} <  - {2x}$ , 解得 $x < \dfrac{1}{6}$ 或 $x > \dfrac{1}{2}$
    \end{jieda}
    \item 不等式 $\left| {{x}^{2} - 4}\right|  < x + 2$ 的解集为\blkda[4]{$(1, 3)$}. 
    \begin{jieda}
        $\left| {{x}^{2} - 4}\right|  < x + 2 \Longleftrightarrow -(x+2) < x^2-4 < (x+2) \Longleftrightarrow \left\{\begin{array}{l}
             (x+2)(x-3) < 0  \\
             (x+2)(x-1) > 0
        \end{array}\right.$，解得$ x \in (1,3)$
    \end{jieda}

    \item 不等式 $\left| {{2x} + 1}\right|  - \left| {5 - x}\right|  > 0$ 的解集为\blkda[4]{$\left( {-\infty , - 6}\right)  \cup  \left( {\dfrac{4}{3}, + \infty }\right)$}
    \begin{jieda}
        由不等式 $\left| {{2x} + 1}\right|  - \left| {5 - x}\right|  > 0$，可得 ${\left( 2x + 1\right) }^{2} > {\left( 5 - x\right) }^{2}$，即 $\left( {x + 6}\right) \left( {{3x} - 4}\right)  > 0$
        
        解得 $x \in  \left( {-\infty , - 6}\right)  \cup  \left( {\dfrac{4}{3}, + \infty }\right)$ .
    \end{jieda}
\end{enumerate}

\subsection*{根据单调性解不等式}
\begin{enumerate}
    \item (2024$\cdot$静安区一模$\cdot$10)不等式 ${\log }_{2}x + \dfrac{x}{2} < 4$ 的解集为\blkda[2]{$\left( {0,4}\right)$}.
    
    \item (2025$\cdot$长宁区一模$\cdot$9)已知 $\alpha  : {2}^{x} + {\log }_{2}x \leq  2,\beta  : x < m$ ，若 $\alpha$ 是 $\beta$ 的充分条件，则实数 $m$ 的取值范围是\blkda[2]{$(1, +\infty)$}.
    \begin{jieda}
        $y = {2}^{x} + {\log }_{2}x$严增，$\alpha : x \in (-\infty, 1]$，$(-\infty, 1] \subseteq (-\infty, m)$，故$m > 1$.
    \end{jieda}

    \item 不等式 ${27}^{x} + 7{\log }_{5}\left( {{36x} + 1}\right)  < {23}$ 的解集为\blkda{$\left( {-\dfrac{1}{36},\dfrac{2}{3}}\right)$}.
\end{enumerate}